\documentclass[11pt]{article}
% Shared preamble for per-Python-file documentation (input via \input{...}).
\usepackage[margin=1in]{geometry}
\usepackage[T1]{fontenc}
\usepackage[utf8]{inputenc}
\usepackage{lmodern}
\usepackage{hyperref}
\usepackage{xcolor}
\usepackage{listings}
\usepackage{listingsutf8}
\lstset{
  basicstyle=\ttfamily\small,
  columns=fullflexible,
  breaklines=true,
  upquote=true,
  inputencoding=utf8,
  extendedchars=true,
  frame=single,
  framerule=0.2pt,
  tabsize=2
}


\title{\texttt{model/wd.py} (Q\&A)}
\author{}
\date{}

\begin{document}
\maketitle

\paragraph{Q: What is the purpose of \texttt{model/wd.py}?}
\textbf{A:} It implements a Wide \& Deep baseline model used for watch-time prediction. The architecture combines an MLP over embedded features (deep part) with optional linear contributions from continuous features (wide part).

\paragraph{Q: What are the inputs and outputs of the model?}
\textbf{A:} The model takes a feature dictionary \lstinline|x_dict| and outputs a sigmoid-activated scalar per example (in this repo, training scripts may rescale labels/predictions, e.g. multiplying labels by 20 in \lstinline|run_vr.py|).

\paragraph{Q: How are features encoded?}
\textbf{A:} Sparse and sequence features are embedded; sequence embeddings are pooled using a mask; continuous features contribute via a learned \lstinline|Linear(1,1)| per feature. The embedded features are concatenated and passed through an MLP.

\paragraph{Q: How does the model combine wide and deep parts?}
\textbf{A:} The deep output is the MLP output; if continuous linear terms exist, they are summed and added to the deep output before applying the final sigmoid.

\paragraph{Q: Code example: the wide+deep combination?}
\textbf{A:}
\begin{lstlisting}[language=Python]
res = self.mlp(emb)
if len(linears) > 0:
    linear_part = torch.concat(linears, dim=1).sum(dim=1, keepdims=True)
    res += linear_part
return torch.sigmoid(res.squeeze(1))
\end{lstlisting}

\paragraph{Q: Full source code?}
\textbf{A:}
\lstinputlisting[language=Python]{/workspace/model/wd.py}

\end{document}
